%% ============================================================
%%  Evaluador de Platillos — Reporte de Avances Intermedios
%%  Universidad de Guanajuato
%%
%%  INSTRUCCIONES PARA OVERLEAF:
%%  Sube las siguientes imágenes a una carpeta llamada "imagenes/"
%%  dentro del proyecto de Overleaf:
%%
%%    imagenes/raw/   → CDMX_Limon_Lateral_Chef_0.jpg … _9.jpg
%%    imagenes/seg/   → CDMX_Limon_Lateral_Chef_0_seg.png … _9_seg.png
%%    imagenes/mask/  → CDMX_Limon_Lateral_Chef_0_mask.png … _9_mask.png
%%    imagenes/pre/   → CDMX_Limon_Lateral_Chef_0_pre.png … _9_pre.png
%%
%%  Compile con: pdflatex → bibtex → pdflatex × 2
%%  (o simplemente pdflatex si no hay referencias externas)
%% ============================================================

\documentclass[12pt,a4paper]{article}

\usepackage[utf8]{inputenc}
\usepackage[T1]{fontenc}
\usepackage[spanish,mexico]{babel}
\usepackage{geometry}
\geometry{margin=2.5cm}

\usepackage{graphicx}
\usepackage{subcaption}
\usepackage{float}
\usepackage{booktabs}
\usepackage{array}
\usepackage{multirow}
\usepackage{tabularx}
\usepackage{longtable}
\usepackage{amsmath}
\usepackage{amssymb}
\usepackage{xcolor}
\usepackage{hyperref}
\usepackage{caption}
\usepackage{enumitem}
\usepackage{parskip}

\usepackage{tikz}
\usetikzlibrary{shapes.geometric, arrows.meta, positioning, fit, calc, backgrounds}

\graphicspath{{imagenes/}}

%% ── Color palette ────────────────────────────────────────────
\definecolor{ugBlue}{RGB}{0,63,135}
\definecolor{ugGold}{RGB}{200,160,0}
\definecolor{pipeGray}{RGB}{220,225,230}
\definecolor{pipeBlue}{RGB}{210,225,245}
\definecolor{pipeGreen}{RGB}{210,240,215}
\definecolor{pipeOrange}{RGB}{255,235,210}

%% ── TikZ styles ──────────────────────────────────────────────
\tikzset{
  block/.style  = {rectangle, rounded corners=4pt, draw=ugBlue, thick,
                   fill=pipeBlue, text width=4.2cm, align=center,
                   minimum height=1.0cm, font=\small},
  blockG/.style = {block, fill=pipeGreen},
  blockO/.style = {block, fill=pipeOrange},
  io/.style     = {trapezium, trapezium left angle=75, trapezium right angle=105,
                   draw=ugBlue, thick, fill=ugGold!30,
                   text width=3.8cm, align=center,
                   minimum height=0.9cm, font=\small},
  feat/.style   = {rectangle, rounded corners=3pt, draw=ugBlue!60, thick,
                   fill=pipeGray, text width=2.6cm, align=center,
                   minimum height=0.85cm, font=\footnotesize},
  arrow/.style  = {-{Stealth[length=6pt]}, ugBlue, thick},
  arrowG/.style = {arrow, ugBlue!40},
}

%% ── Metadata ─────────────────────────────────────────────────
\title{%
  \textcolor{ugBlue}{\textbf{Evaluador Automático de Platillos}}\\[4pt]
  \large Avances Intermedios: Preprocesamiento de Imágenes\\
  y Extracción de Características Visuales
}
\author{%
  Universidad de Guanajuato\\
  Procesamiento de Imágenes y Visión por Computadora
}
\date{\today}

%% ============================================================
\begin{document}
\maketitle
\tableofcontents
\vspace{1em}
\hrule
\vspace{1em}

%% ============================================================
\section{Introducción}
%% ============================================================

El proyecto \textit{Evaluador de Platillos} tiene como objetivo
construir un sistema de evaluación automática de la calidad del
emplatado culinario.  Dado el plato de un \textbf{Chef} (referencia)
y el plato de un \textbf{Estudiante}, el sistema debe estimar,
mediante descriptores de visión por computadora y un sistema difuso,
qué tan similar es la presentación del Estudiante respecto al estándar
del Chef.

Este reporte describe los avances obtenidos en las dos primeras etapas
del pipeline:
\begin{enumerate}[nosep]
  \item \textbf{Preprocesamiento de imágenes}: segmentación automática
        del plato (GrabCut), limpieza morfológica, inpainting, recorte
        al contenido y normalización a 224$\times$224 px.
  \item \textbf{Extracción de características}: seis descriptores
        basados en ventanas deslizantes (RAW, SDH, GLCM, LBP, HOG y
        Color).
\end{enumerate}

Las pruebas se realizaron sobre un conjunto de \textbf{10 imágenes}
del platillo \textit{Limón} en vista lateral, tomadas por el mismo
Chef desde distintos ángulos de cámara (CDMX, resolución 3840×2160
px).

%% ============================================================
\section{Flujo de Trabajo del Pipeline}
%% ============================================================

La figura~\ref{fig:pipeline} muestra el diagrama completo del
pipeline implementado en C++17 con OpenCV~4.6.

\begin{figure}[H]
\centering
\begin{tikzpicture}[node distance=0.55cm and 1.8cm]

  %% ── Entrada ───────────────────────────────────────────────
  \node[io] (entrada) {Imagen Original\\(JPEG 3840×2160)};

  %% ── Etapa 1: Segmentación ─────────────────────────────────
  \node[block, below=of entrada] (grabcut)
    {GrabCut\\ROI central 70\%\\5 iteraciones};
  \node[block, below=of grabcut] (morfo)
    {Postprocesamiento Morfológico\\Opening 5×5 ×2\\Closing 15×15 ×3};
  \node[block, below=of morfo] (inpaint)
    {Inpainting (Telea)\\Relleno de huecos};

  %% ── Etapa 2: Preprocesamiento ─────────────────────────────
  \node[blockG, below=of inpaint] (crop)
    {Recorte al Contenido\\Bounding box + 5\% padding};
  \node[blockG, below=of crop] (resize)
    {Redimensionado\\224 × 224 px};
  \node[blockG, below=of resize] (orient)
    {Corrección de Orientación\\minAreaRect, límite ±30°};
  \node[blockG, below=of orient] (ssim)
    {Alineación SSIM\\vs.\ referencia Chef};

  %% ── Salida preprocesada ───────────────────────────────────
  \node[io, below=of ssim] (prepro)
    {Imagen Preprocesada\\(PNG 224×224)};

  %% ── Extracción de características ─────────────────────────
  \node[blockO, below=of prepro] (feat_hub)
    {Extracción de\\Características\\(ventanas deslizantes)};

  %% Seis descriptores en dos filas
  \node[feat, below left=0.9cm and 3.4cm of feat_hub]  (raw)   {RAW\\(168 dim.)};
  \node[feat, below left=0.9cm and 0.8cm of feat_hub]  (sdh)   {SDH\\(168 dim.)};
  \node[feat, below=0.9cm       of feat_hub]            (glcm)  {GLCM\\(2\,016 dim.)};
  \node[feat, below right=0.9cm and 0.8cm of feat_hub] (lbp)   {LBP\\(120 dim.)};
  \node[feat, below right=0.9cm and 3.4cm of feat_hub] (hog)   {HOG\\(26\,244 dim.)};
  \node[feat, below right=0.9cm and 6.0cm of feat_hub] (color) {Color\\(648 dim.)};

  %% CSV
  \node[io, below=2.6cm of feat_hub] (csv)
    {Vectores de Características\\(archivos CSV)};

  %% ── Flechas ───────────────────────────────────────────────
  \draw[arrow] (entrada) -- (grabcut);
  \draw[arrow] (grabcut) -- (morfo);
  \draw[arrow] (morfo)   -- (inpaint);
  \draw[arrow] (inpaint) -- (crop);
  \draw[arrow] (crop)    -- (resize);
  \draw[arrow] (resize)  -- (orient);
  \draw[arrow] (orient)  -- (ssim);
  \draw[arrow] (ssim)    -- (prepro);
  \draw[arrow] (prepro)  -- (feat_hub);

  \draw[arrowG] (feat_hub.south) -- (raw.north);
  \draw[arrowG] (feat_hub.south) -- (sdh.north);
  \draw[arrowG] (feat_hub.south) -- (glcm.north);
  \draw[arrowG] (feat_hub.south) -- (lbp.north);
  \draw[arrowG] (feat_hub.south) -- (hog.north);
  \draw[arrowG] (feat_hub.south) -- (color.north);

  \draw[arrowG] (raw.south)   |- (csv.west);
  \draw[arrowG] (sdh.south)   |- (csv.west);
  \draw[arrowG] (glcm.south)  -- (csv.north);
  \draw[arrowG] (lbp.south)   |- (csv.east);
  \draw[arrowG] (hog.south)   |- (csv.east);
  \draw[arrowG] (color.south) |- (csv.east);

  %% ── Etiquetas de etapa ────────────────────────────────────
  \begin{scope}[on background layer]
    \node[draw=ugBlue!30, dashed, rounded corners,
          fill=ugBlue!5, fit=(grabcut)(inpaint),
          label={[ugBlue,font=\scriptsize]left:Etapa 1\\Segmentación}]
          {};
    \node[draw=ugGold!60, dashed, rounded corners,
          fill=ugGold!5, fit=(crop)(ssim),
          label={[ugGold!80!black,font=\scriptsize]left:Etapa 2\\Preproceso}]
          {};
  \end{scope}

\end{tikzpicture}
\caption{Pipeline completo de preprocesamiento y extracción de
         características del sistema evaluador de platillos.}
\label{fig:pipeline}
\end{figure}

%% ============================================================
\section{Implementación del Preprocesamiento}
%% ============================================================

\subsection{Segmentación con GrabCut}

GrabCut~\cite{rother2004grabcut} es un algoritmo iterativo de
segmentación de primer plano basado en Grafos y Mezclas Gaussianas
(GMM).  Se utilizó la siguiente configuración:

\begin{itemize}[nosep]
  \item \textbf{ROI automática}: rectángulo central del 70\% de la
        imagen (margen del 15\% por lado), evitando bordes de mesa.
  \item \textbf{Iteraciones}: 5 ciclos de refinamiento EM-GMM.
  \item \textbf{Salida}: máscara binaria de primer plano.
\end{itemize}

\subsection{Postprocesamiento Morfológico}

Para eliminar huecos y ruido en la máscara de GrabCut se aplicaron
dos operaciones morfológicas con elementos estructurantes elípticos:

\begin{enumerate}[nosep]
  \item \textbf{Opening} (erode + dilate): kernel elíptico 5×5,
        2~repeticiones.  Elimina pequeñas regiones falsas de primer
        plano.
  \item \textbf{Closing} (dilate + erode): kernel elíptico 15×15,
        3~repeticiones.  Rellena huecos internos del plato.
\end{enumerate}

\subsection{Inpainting}

Los píxeles clasificados como fondo en la máscara pero rodeados de
primer plano (agujeros) se rellenan con el método de Telea
(\texttt{cv::inpaint}, radio 3~px) para evitar artefactos en la
imagen segmentada.

\subsection{Recorte al Contenido y Redimensionado}

Tras la segmentación se calcula el bounding box de los píxeles no
nulos, se agrega un margen del 5\% y se recorta la imagen.  A
continuación se redimensiona al tamaño estándar de
\textbf{224$\times$224~px} mediante interpolación cúbica, colocando el
plato sobre fondo negro.

\subsection{Corrección de Orientación y Alineación}

\paragraph{Corrección de orientación.}
Se calcula el rectángulo de área mínima del mayor contorno de la
máscara (\texttt{cv::minAreaRect}).  Si el ángulo de rotación
detectado está dentro del rango $|\theta| \leq 30°$ se aplica una
rotación compensatoria; en caso contrario la imagen se acepta sin
cambios (valores mayores indican ruido en la máscara, no una
inclinación real).

El canvas expandido para la rotación utiliza:
\[
  w' = \lceil W|\cos\theta| + H|\sin\theta| \rceil, \quad
  h' = \lceil W|\sin\theta| + H|\cos\theta| \rceil
\]
con valores absolutos para manejar correctamente ángulos obtusos.

\paragraph{Alineación SSIM.}
Cuando existen múltiples imágenes del mismo platillo, la primera
imagen del Chef se usa como referencia.  Cada imagen subsecuente se
compara con la referencia en orientación original y rotada 180°;
se selecciona la orientación con mayor SSIM
(\textit{Structural Similarity Index}).

%% ============================================================
\section{Resultados del Preprocesamiento}
%% ============================================================

\subsection{Ejemplos Visuales}

Las figuras~\ref{fig:ejemplo0}--\ref{fig:ejemplo2} muestran el
proceso completo para tres imágenes representativas del conjunto de
datos.

%% ── Figura 1 ─────────────────────────────────────────────────
\begin{figure}[H]
\centering
\begin{subfigure}[t]{0.24\textwidth}
  \centering
  \includegraphics[width=\textwidth]{raw/CDMX_Limon_Lateral_Chef_0.jpg}
  \caption{Original (3840×2160)}
\end{subfigure}\hfill
\begin{subfigure}[t]{0.24\textwidth}
  \centering
  \includegraphics[width=\textwidth]{mask/CDMX_Limon_Lateral_Chef_0_mask.png}
  \caption{Máscara GrabCut}
\end{subfigure}\hfill
\begin{subfigure}[t]{0.24\textwidth}
  \centering
  \includegraphics[width=\textwidth]{seg/CDMX_Limon_Lateral_Chef_0_seg.png}
  \caption{Segmentada + Inpainting}
\end{subfigure}\hfill
\begin{subfigure}[t]{0.24\textwidth}
  \centering
  \includegraphics[width=\textwidth]{pre/CDMX_Limon_Lateral_Chef_0_pre.png}
  \caption{Preprocesada 224×224}
\end{subfigure}
\caption{Imagen \texttt{Chef\_0}: etapas del preprocesamiento.}
\label{fig:ejemplo0}
\end{figure}

%% ── Figura 2 ─────────────────────────────────────────────────
\begin{figure}[H]
\centering
\begin{subfigure}[t]{0.24\textwidth}
  \centering
  \includegraphics[width=\textwidth]{raw/CDMX_Limon_Lateral_Chef_3.jpg}
  \caption{Original}
\end{subfigure}\hfill
\begin{subfigure}[t]{0.24\textwidth}
  \centering
  \includegraphics[width=\textwidth]{mask/CDMX_Limon_Lateral_Chef_3_mask.png}
  \caption{Máscara GrabCut}
\end{subfigure}\hfill
\begin{subfigure}[t]{0.24\textwidth}
  \centering
  \includegraphics[width=\textwidth]{seg/CDMX_Limon_Lateral_Chef_3_seg.png}
  \caption{Segmentada}
\end{subfigure}\hfill
\begin{subfigure}[t]{0.24\textwidth}
  \centering
  \includegraphics[width=\textwidth]{pre/CDMX_Limon_Lateral_Chef_3_pre.png}
  \caption{Preprocesada}
\end{subfigure}
\caption{Imagen \texttt{Chef\_3}: etapas del preprocesamiento.}
\label{fig:ejemplo1}
\end{figure}

%% ── Figura 3 ─────────────────────────────────────────────────
\begin{figure}[H]
\centering
\begin{subfigure}[t]{0.24\textwidth}
  \centering
  \includegraphics[width=\textwidth]{raw/CDMX_Limon_Lateral_Chef_7.jpg}
  \caption{Original}
\end{subfigure}\hfill
\begin{subfigure}[t]{0.24\textwidth}
  \centering
  \includegraphics[width=\textwidth]{mask/CDMX_Limon_Lateral_Chef_7_mask.png}
  \caption{Máscara GrabCut}
\end{subfigure}\hfill
\begin{subfigure}[t]{0.24\textwidth}
  \centering
  \includegraphics[width=\textwidth]{seg/CDMX_Limon_Lateral_Chef_7_seg.png}
  \caption{Segmentada}
\end{subfigure}\hfill
\begin{subfigure}[t]{0.24\textwidth}
  \centering
  \includegraphics[width=\textwidth]{pre/CDMX_Limon_Lateral_Chef_7_pre.png}
  \caption{Preprocesada}
\end{subfigure}
\caption{Imagen \texttt{Chef\_7}: etapas del preprocesamiento.}
\label{fig:ejemplo2}
\end{figure}

\subsection{Resumen del Procesamiento}

\begin{table}[H]
\centering
\caption{Resultados del pipeline de preprocesamiento para las 10
         imágenes del conjunto de prueba.}
\label{tab:prepro_resumen}
\begin{tabular}{@{}llccc@{}}
\toprule
\textbf{Imagen} & \textbf{Ciudad/Platillo} & \textbf{Res.\ orig.} &
  \textbf{Segmentada} & \textbf{Preprocesada} \\
\midrule
Chef\_0 & CDMX / Limón Lateral & 3840×2160 & \checkmark & 224×224 \\
Chef\_1 & CDMX / Limón Lateral & 3840×2160 & \checkmark & 224×224 \\
Chef\_2 & CDMX / Limón Lateral & 3840×2160 & \checkmark & 224×224 \\
Chef\_3 & CDMX / Limón Lateral & 3840×2160 & \checkmark & 224×224 \\
Chef\_4 & CDMX / Limón Lateral & 3840×2160 & \checkmark & 224×224 \\
Chef\_5 & CDMX / Limón Lateral & 3840×2160 & \checkmark & 224×224 \\
Chef\_6 & CDMX / Limón Lateral & 3840×2160 & \checkmark & 224×224 \\
Chef\_7 & CDMX / Limón Lateral & 3840×2160 & \checkmark & 224×224 \\
Chef\_8 & CDMX / Limón Lateral & 3840×2160 & \checkmark & 224×224 \\
Chef\_9 & CDMX / Limón Lateral & 3840×2160 & \checkmark & 224×224 \\
\midrule
\multicolumn{2}{@{}l}{\textbf{Total}} & --- & \textbf{10/10} & \textbf{10/10} \\
\bottomrule
\end{tabular}
\end{table}

%% ============================================================
\section{Extractores de Características}
%% ============================================================

\subsection{Descripción General}

Todos los descriptores (excepto HOG) se calculan mediante una
\textbf{ventana deslizante} de paso 1~px con doce tamaños:
$w \in \{3, 5, 7, 9, 11, 13, 15, 17, 19, 21, 23, 25\}$.
Para cada posición de ventana se computan las características
internas; luego se calcula la \textbf{media} ($\mu$) y la
\textbf{desviación estándar} ($\sigma$) del mapa resultante.
El vector final por imagen concatena $(\mu, \sigma)$ de todos los
tamaños y características.

\begin{table}[H]
\centering
\caption{Dimensionalidad de los vectores de características generados
         por cada extractor.}
\label{tab:dimensiones}
\begin{tabular}{@{}lrrrc@{}}
\toprule
\textbf{Extractor} & \textbf{Ventanas} & \textbf{Feats/ventana} &
  \textbf{Estad.} & \textbf{Total (sin IDs)} \\
\midrule
RAW   & 12 & 7  & $\mu, \sigma$ & 168  \\
SDH   & 12 & 7  & $\mu, \sigma$ & 168  \\
GLCM  & 12 & 84 & $\mu, \sigma$ & 2\,016 \\
LBP   & 12 & 5  & $\mu, \sigma$ & 120  \\
HOG   & —  & — & —             & 26\,244 \\
Color & 12 & 27 & $\mu, \sigma$ & 648  \\
\midrule
\multicolumn{4}{@{}l}{\textbf{Total acumulado}} & \textbf{29\,368} \\
\bottomrule
\end{tabular}
\end{table}

\subsection{Descriptor RAW (Estadísticas de Histograma)}

El extractor \textit{Raw} calcula siete estadísticos de la
distribución de intensidades en escala de grises de cada ventana,
análogos a las características de co-ocurrencia de Haralick pero
aplicados directamente al histograma de píxeles:
Energía, Contraste, Correlación, Homogeneidad,
IDF (\textit{Inverse Difference Feature}), Entropía y Varianza.

La tabla~\ref{tab:raw_w3} muestra los valores medios de cada
característica para ventana $w=3$ en las 10 imágenes del conjunto.

\begin{table}[H]
\centering
\caption{Características RAW ($w=3$, estadístico $\mu$) para las
         10 imágenes preprocesadas de Chef/Limón/Lateral/CDMX.}
\label{tab:raw_w3}
\setlength{\tabcolsep}{5pt}
\begin{tabular}{@{}lccccccc@{}}
\toprule
\textbf{Imagen} & \textbf{Energía} & \textbf{Contraste} &
  \textbf{Correlación} & \textbf{Homog.} & \textbf{IDF} &
  \textbf{Entropía} & \textbf{Varianza} \\
\midrule
Chef\_0 & 0.6231 & 169.07 & 0.5849 & 0.7095 & 0.7155 & 1.0696 & 169.07 \\
Chef\_1 & 0.6434 & 164.37 & 0.5457 & 0.7195 & 0.7213 & 1.0276 & 164.37 \\
Chef\_2 & 0.6235 & 176.79 & 0.5759 & 0.7066 & 0.7097 & 1.0782 & 176.79 \\
Chef\_3 & 0.6213 & 159.24 & 0.5896 & 0.7095 & 0.7159 & 1.0727 & 159.24 \\
Chef\_4 & 0.6253 & 164.00 & 0.5673 & 0.7098 & 0.7117 & 1.0731 & 164.00 \\
Chef\_5 & 0.6135 & 167.41 & 0.6067 & 0.7014 & 0.7087 & 1.0935 & 167.41 \\
Chef\_6 & 0.5936 & 164.98 & 0.6402 & 0.6942 & 0.7047 & 1.1379 & 164.98 \\
Chef\_7 & 0.6383 & 165.00 & 0.5620 & 0.7182 & 0.7229 & 1.0302 & 165.00 \\
Chef\_8 & 0.6099 & 167.86 & 0.6104 & 0.6977 & 0.7045 & 1.1059 & 167.86 \\
Chef\_9 & 0.6093 & 171.81 & 0.5942 & 0.6982 & 0.7011 & 1.1175 & 171.81 \\
\midrule
\textbf{Media} & \textbf{0.6201} & \textbf{167.05} & \textbf{0.5877} &
  \textbf{0.7065} & \textbf{0.7116} & \textbf{1.0806} & \textbf{167.05} \\
\textbf{Desv.} & 0.0151 & 4.97 & 0.0305 & 0.0082 & 0.0083 & 0.0326 & 4.97 \\
\bottomrule
\end{tabular}
\end{table}

\noindent\textbf{Observación}: Los valores de Contraste y Varianza
coinciden numéricamente porque ambas métricas miden la dispersión
cuadrática de intensidades en el espacio de histograma RAW; se
conservan como características distintas para mantener consistencia
con la base de características Haralick.

\subsection{Descriptor SDH (Sum-Difference Histograms)}

El extractor SDH calcula el histograma de suma ($s = I_i + I_j$) y
de diferencia ($d = |I_i - I_j|$) de pares de píxeles vecinos,
con 64 niveles de cuantización.  Sobre los histogramas concatenados
se computan los mismos siete estadísticos que en RAW.

\begin{table}[H]
\centering
\caption{Características SDH ($w=3$, estadístico $\mu$) — valores
         representativos.}
\label{tab:sdh_w3}
\setlength{\tabcolsep}{5pt}
\begin{tabular}{@{}lcccc@{}}
\toprule
\textbf{Imagen} & \textbf{Media} & \textbf{Varianza} & \textbf{Correlación} & \textbf{Correlación $\sigma$} \\
\midrule
Chef\_0 & 104.39 & 108.94 & $-$0.113 & 0.356 \\
Chef\_1 &  98.12 & 105.87 & $-$0.101 & 0.357 \\
Chef\_2 & 103.22 & 113.78 & $-$0.110 & 0.360 \\
Chef\_3 & 106.22 & 102.68 & $-$0.124 & 0.359 \\
Chef\_4 & 100.07 & 105.73 & $-$0.105 & 0.351 \\
Chef\_5 & 110.04 & 107.90 & $-$0.123 & 0.363 \\
Chef\_6 & 110.85 & 106.40 & $-$0.135 & 0.368 \\
Chef\_7 & 101.26 & 106.26 & $-$0.112 & 0.358 \\
Chef\_8 & 110.96 & 108.46 & $-$0.127 & 0.367 \\
Chef\_9 & 104.44 & 110.80 & $-$0.107 & 0.358 \\
\midrule
\textbf{Media} & 104.96 & 107.68 & $-$0.116 & 0.360 \\
\bottomrule
\end{tabular}
\end{table}

\noindent La Correlación SDH negativa refleja que píxeles vecinos
tienden a diferir (texturas con transición fondo negro/plato), y
permanece estable entre imágenes ($\sigma \approx 0.36$).

\subsection{Descriptor GLCM (Características de Haralick)}

La Matriz de Co-ocurrencia de Niveles de Grises (GLCM) se calcula
para 3 distancias ($d \in \{1,3,7\}$) y 4 ángulos
($\theta \in \{0°, 45°, 90°, 135°\}$), generando $3 \times 4 = 12$
configuraciones.  Para cada configuración se extraen 7 características
de Haralick: Energía, Contraste, Correlación, Homogeneidad, IDF,
Entropía y Varianza.  Total: $12 \times 7 = 84$ características por
ventana.

La GLCM se normaliza a 64 niveles (\texttt{GLCM\_LEVELS=64}) mediante:
\[
  q[r,c] = \left\lfloor I[r,c] \cdot \frac{63}{255} \right\rfloor
\]

La tabla~\ref{tab:glcm_w3} muestra la Energía para $w=3$ en las 4
orientaciones de la configuración $d=1$.

\begin{table}[H]
\centering
\caption{Energía GLCM ($w=3$, $d=1$) en las cuatro orientaciones.}
\label{tab:glcm_w3}
\setlength{\tabcolsep}{5pt}
\begin{tabular}{@{}lcccc@{}}
\toprule
\textbf{Imagen} & $\theta=0°$ & $\theta=45°$ & $\theta=90°$ & $\theta=135°$ \\
\midrule
Chef\_0 & 0.7252 & 0.7396 & 0.7332 & 0.7399 \\
Chef\_1 & 0.7310 & 0.7465 & 0.7394 & 0.7454 \\
Chef\_2 & 0.7157 & 0.7323 & 0.7245 & 0.7314 \\
Chef\_3 & 0.7188 & 0.7356 & 0.7274 & 0.7331 \\
Chef\_4 & 0.7142 & 0.7295 & 0.7232 & 0.7299 \\
Chef\_5 & 0.7096 & 0.7260 & 0.7191 & 0.7256 \\
Chef\_6 & 0.7086 & 0.7232 & 0.7161 & 0.7246 \\
Chef\_7 & 0.7250 & 0.7393 & 0.7330 & 0.7408 \\
Chef\_8 & 0.7085 & 0.7234 & 0.7179 & 0.7235 \\
Chef\_9 & 0.7024 & 0.7200 & 0.7121 & 0.7194 \\
\midrule
\textbf{Media} & 0.7159 & 0.7315 & 0.7246 & 0.7314 \\
\bottomrule
\end{tabular}
\end{table}

\noindent La Energía GLCM superior en $\theta=45°$ y $135°$ (direcciones
diagonales) respecto a $\theta=0°$ es consistente con la textura
circular del plato/limón, que produce mayor regularidad en co-ocurrencias
diagonales.

\subsection{Descriptor LBP (Local Binary Patterns)}

LBP codifica la textura local comparando cada píxel con sus 8
vecinos en radio 1.  El código binario resultante ($\in[0,255]$) se
procesa en cada ventana para obtener 5 estadísticos:
Media, Varianza, Correlación, Contraste y Homogeneidad.

\begin{table}[H]
\centering
\caption{Características LBP ($w=3$) — Media e índice de Correlación.}
\label{tab:lbp_w3}
\setlength{\tabcolsep}{5pt}
\begin{tabular}{@{}lcccc@{}}
\toprule
\textbf{Imagen} & \textbf{Media$_\mu$} & \textbf{Media$_\sigma$} &
  \textbf{Correlación$_\mu$} & \textbf{Contraste$_\mu$} \\
\midrule
Chef\_0 & 202.16 & 61.68 & 0.6117 & 2781.12 \\
Chef\_1 & 206.28 & 60.89 & 0.5765 & 2510.95 \\
Chef\_2 & 203.35 & 61.68 & 0.6088 & 2609.17 \\
Chef\_3 & 201.36 & 61.73 & 0.6209 & 2849.62 \\
Chef\_4 & 203.32 & 61.21 & 0.5926 & 2608.50 \\
Chef\_5 & 201.51 & 61.52 & 0.6400 & 2823.78 \\
Chef\_6 & 198.52 & 59.78 & 0.6609 & 3194.93 \\
Chef\_7 & 204.65 & 61.60 & 0.5918 & 2546.53 \\
Chef\_8 & 201.47 & 60.84 & 0.6477 & 2890.02 \\
Chef\_9 & 200.15 & 62.99 & 0.6209 & 2767.58 \\
\midrule
\textbf{Media} & 202.28 & 61.39 & 0.6172 & 2758.22 \\
\textbf{Desv.}  &   2.10 &  0.84 & 0.0269 & 194.24  \\
\bottomrule
\end{tabular}
\end{table}

\noindent La Media LBP elevada ($\approx 202$ sobre 255) indica que la
mayoría de los píxeles son más brillantes que sus vecinos (lo que es
consistente con el plato blanco sobre fondo negro), y el Contraste
refleja la heterogeneidad de los patrones binarios locales.

\subsection{Descriptor HOG (Histogram of Oriented Gradients)}

HOG se aplica a la imagen completa 224×224 (sin ventana deslizante):
\begin{itemize}[nosep]
  \item Celdas de 8×8 píxeles → $28 \times 28 = 784$ celdas.
  \item Bloques de 2×2 celdas → $27 \times 27 = 729$ bloques.
  \item 9 orientaciones por celda; 36 valores por bloque.
  \item Total: $729 \times 36 = \mathbf{26\,244}$ dimensiones.
\end{itemize}

\subsection{Descriptor de Color (HSV + CIE L*a*b*)}

Se extraen estadísticos de color en los espacios HSV y CIE~L*a*b*:
media, desviación, sesgo y curtosis para cada canal, más la varianza
cruzada entre canales.  Total: 27 características por ventana.

\begin{table}[H]
\centering
\caption{Características de Color ($w=3$) — medias de canales HSV y L* para las 10 imágenes.}
\label{tab:color_w3}
\setlength{\tabcolsep}{5pt}
\begin{tabular}{@{}lcccc@{}}
\toprule
\textbf{Imagen} & \textbf{H$_\mu$} & \textbf{S$_\mu$} & \textbf{V$_\mu$} & \textbf{L*$_\mu$} \\
\midrule
Chef\_0 & 25.71 & 34.75 & 109.70 & 109.46 \\
Chef\_1 & 25.52 & 36.75 & 103.57 & 103.10 \\
Chef\_2 & 26.82 & 35.95 & 108.48 & 108.39 \\
Chef\_3 & 25.90 & 33.30 & 111.16 & 111.36 \\
Chef\_4 & 24.75 & 35.05 & 105.32 & 105.10 \\
Chef\_5 & 26.68 & 35.05 & 115.32 & 115.38 \\
Chef\_6 & 25.05 & 34.85 & 116.26 & 116.24 \\
Chef\_7 & 23.28 & 34.90 & 106.66 & 106.31 \\
Chef\_8 & 25.33 & 36.99 & 116.66 & 116.43 \\
Chef\_9 & 24.24 & 36.38 & 110.01 & 109.67 \\
\midrule
\textbf{Media} & 25.31 & 35.40 & 110.31 & 110.14 \\
\textbf{Desv.} &  1.00 &  1.16 &   4.49 &   4.52 \\
\bottomrule
\end{tabular}
\end{table}

\noindent El Tono (H) cercano a 25° corresponde al amarillo-naranja
del limón.  La Saturación (S$\approx$35) moderada refleja la mezcla
entre el plato blanco (S≈0) y la pulpa del limón.  L* y V son
casi idénticos, como esperado en regiones poco saturadas.

%% ============================================================
\section{Discusión de Resultados}
%% ============================================================

\subsection{Variabilidad entre Imágenes}

Analizando las tablas de características con $w=3$ se observa:

\begin{itemize}[nosep]
  \item \textbf{RAW}: Energía entre 0.594 (Chef\_6) y 0.643 (Chef\_1).
        Contraste oscila 159–177, sensible al ángulo de iluminación.
        La Correlación (0.546–0.640) muestra el mayor rango relativo.
  \item \textbf{SDH}: Correlación siempre negativa ($\approx -0.12$),
        indicando que pares de píxeles vecinos tienden a diferir
        (transición plato/fondo predominante).
  \item \textbf{GLCM}: Energía $d=1$ entre 0.702 y 0.731, mayor en
        $\theta=45°$ y $135°$ (simetría circular del plato).
  \item \textbf{LBP}: Media $\approx 202$ (escala 0–255), con poca
        varianza entre imágenes ($\sigma=2.1$): la textura local del
        plato blanco es homogénea.
  \item \textbf{Color}: Tono $\approx 25°$ (limón/amarillo).
        Luminosidad L* y V difieren apenas 0.2 unidades.
  \item La \textbf{desviación estándar inter-imagen} de todas las
        métricas es $< 5\%$ del rango, confirmando varianza
        intra-clase baja — deseable para el clasificador difuso.
\end{itemize}

\subsection{Corrección Identificada: Inicialización de Mapas}

Durante la implementación se detectó un error sutil de C++/OpenCV:
\begin{quote}
  \texttt{std::vector<cv::Mat> maps(N, cv::Mat::zeros(...));}
\end{quote}
crea $N$ objetos \texttt{cv::Mat} que comparten el mismo buffer de
datos (copia superficial por el constructor de copia de OpenCV).
Esto hacía que todos los canales del mapa de características
almacenaran únicamente el valor de la \textit{última} característica
escrita, corrompiendo las demás columnas del CSV.

La corrección implementada inicializa cada mapa de forma
independiente:
\begin{verbatim}
std::vector<cv::Mat> maps(N);
for (int k = 0; k < N; ++k)
    maps[k] = cv::Mat::zeros(out_rows, out_cols, CV_64F);
\end{verbatim}

Este bug fue corregido en todos los extractores (RAW, SDH, GLCM,
LBP y Color) y los archivos CSV fueron regenerados.

%% ============================================================
\section{Estructura de Datos}
%% ============================================================

\subsection{Jerarquía de Archivos}

\begin{verbatim}
Data/
├── Raw/Imagenes/
│   └── CDMX/Limon/Lateral/Chef/   ← 10 imágenes JPEG 4K
├── Segmentadas/
│   └── CDMX/Limon/Lateral/Chef/   ← *_mask.png  +  *_seg.png
├── Preprocesadas/
│   └── CDMX/Limon/Lateral/Chef/   ← *_pre.png (224×224)
└── Features/
    ├── features_raw.csv   (170 cols)
    ├── features_sdh.csv   (170 cols)
    ├── features_glcm.csv  (2018 cols)
    ├── features_lbp.csv   (122 cols)
    ├── features_hog.csv   (26248 cols)
    └── features_color.csv (650 cols)
\end{verbatim}

\subsection{Convención de Nombres}

Todos los archivos siguen el esquema:
\[
  \texttt{<Ciudad>\_<Platillo>\_<Ángulo>\_<Autor>[\_<N>][\_sufijo].ext}
\]
Ejemplo: \texttt{CDMX\_Limon\_Lateral\_Chef\_3\_pre.png}.

\subsection{Tamaños de los Archivos de Características}

\begin{table}[H]
\centering
\caption{Resumen de archivos CSV generados (10 imágenes de prueba).}
\label{tab:csvs}
\begin{tabular}{@{}lrrr@{}}
\toprule
\textbf{Archivo} & \textbf{Columnas} & \textbf{Filas} & \textbf{Tamaño} \\
\midrule
\texttt{features\_raw.csv}   & 170    & 10 &  18 KB \\
\texttt{features\_sdh.csv}   & 170    & 10 &  19 KB \\
\texttt{features\_glcm.csv}  & 2\,018 & 10 & 222 KB \\
\texttt{features\_lbp.csv}   & 122    & 10 &  13 KB \\
\texttt{features\_hog.csv}   & 26\,248& 10 & 2.0 MB \\
\texttt{features\_color.csv} & 650    & 10 &  67 KB \\
\midrule
\textbf{Total acumulado} & \textbf{29\,378} & --- & \textbf{2.3 MB} \\
\bottomrule
\end{tabular}
\end{table}

\subsection{Formato de los CSV}

Cada fila del CSV corresponde a una imagen:
\[
  \underbrace{\text{label},\,\text{filename}}_{\text{2 cols de id}},\;
  \underbrace{\text{Feat\_w3\_X\_mean},\,\text{Feat\_w3\_X\_std},\;\ldots}_{\text{descriptores}}
\]
La etiqueta (\texttt{label}) codifica la clase como
\texttt{<Platillo>\_<Autor>} (e.g., \texttt{Limon\_Chef}).

%% ============================================================
\section{Conclusiones y Trabajo Futuro}
%% ============================================================

\subsection{Logros de la Etapa Actual}

\begin{enumerate}[nosep]
  \item Pipeline de preprocesamiento completamente automático:
        10/10 imágenes procesadas en un único script de línea de
        comandos.
  \item Seis extractores de características implementados y
        corregidos, generando vectores válidos y reproducibles.
  \item Imágenes normalizadas a 224×224~px con el plato centrado
        sobre fondo negro, coherentes con el \textit{ReporteFinal}
        del proyecto.
\end{enumerate}

\subsection{Trabajo Pendiente}

\begin{itemize}[nosep]
  \item \textbf{Ampliar el conjunto de datos} con imágenes de
        Estudiantes y de otros platillos y ciudades.
  \item \textbf{Análisis de relevancia} de características: reducción
        de dimensionalidad (PCA, selección de características) dada
        la alta dimensionalidad acumulada (29\,368 dims).
  \item \textbf{Diseño del sistema difuso} de evaluación (módulo de
        inferencia Mamdani o Takagi-Sugeno) para mapear los
        descriptores de distancia Chef/Estudiante a una puntuación de
        calidad.
  \item \textbf{Validación} del segmentador GrabCut en platillos con
        mayor variedad cromática (sopas, guisos).
\end{itemize}

%% ============================================================
%% Bibliografía básica (sin BibTeX — usa thebibliography)
%% ============================================================
\begin{thebibliography}{9}

\bibitem{rother2004grabcut}
  C.~Rother, V.~Kolmogorov y A.~Blake,
  ``GrabCut: Interactive Foreground Extraction using Iterated Graph Cuts,''
  \textit{ACM SIGGRAPH}, 2004.

\bibitem{haralick1973textural}
  R.~M.~Haralick, K.~Shanmugam e I.~Dinstein,
  ``Textural Features for Image Classification,''
  \textit{IEEE Trans.\ Systems, Man, and Cybernetics}, 3(6):610--621, 1973.

\bibitem{dalal2005hog}
  N.~Dalal y B.~Triggs,
  ``Histograms of Oriented Gradients for Human Detection,''
  \textit{CVPR}, 2005.

\bibitem{ojala2002lbp}
  T.~Ojala, M.~Pietikäinen y T.~Mäenpää,
  ``Multiresolution Gray-Scale and Rotation Invariant Texture Classification
  with Local Binary Patterns,''
  \textit{IEEE TPAMI}, 24(7):971--987, 2002.

\bibitem{opencv}
  G.~Bradski, ``The OpenCV Library,''
  \textit{Dr.\ Dobb's Journal of Software Tools}, 2000.

\end{thebibliography}

\end{document}
